\documentclass{../base/canon}

%% canon_variables.tex — Valores por defecto de metadatos institucionales
%%
%% Incluir ANTES del \begin{document} para sobreescribir los defaults de la clase.
%% El template puede sobreescribir cualquier valor después de este \input.
%%
%% Uso en template:
%%   %% canon_variables.tex — Valores por defecto de metadatos institucionales
%%
%% Incluir ANTES del \begin{document} para sobreescribir los defaults de la clase.
%% El template puede sobreescribir cualquier valor después de este \input.
%%
%% Uso en template:
%%   %% canon_variables.tex — Valores por defecto de metadatos institucionales
%%
%% Incluir ANTES del \begin{document} para sobreescribir los defaults de la clase.
%% El template puede sobreescribir cualquier valor después de este \input.
%%
%% Uso en template:
%%   \input{../base/canon_variables}
%%   \SetDocTitle{Mi Informe Específico}   % sobreescribe el default

\SetDocLogo{../assets/logo.pdf}
\SetDocUnit{Tecnología \& Sistemas}
\SetContactUnit{Mesa de soporte}
\SetContactMembers{%
  Equipo CoreX & Soporte & soporte@empresa.com%
}
\SetContactNote{Para soporte técnico o consultas sobre este documento.}

%%   \SetDocTitle{Mi Informe Específico}   % sobreescribe el default

\SetDocLogo{../assets/logo.pdf}
\SetDocUnit{Tecnología \& Sistemas}
\SetContactUnit{Mesa de soporte}
\SetContactMembers{%
  Equipo CoreX & Soporte & soporte@empresa.com%
}
\SetContactNote{Para soporte técnico o consultas sobre este documento.}

%%   \SetDocTitle{Mi Informe Específico}   % sobreescribe el default

\SetDocLogo{../assets/logo.pdf}
\SetDocUnit{Tecnología \& Sistemas}
\SetContactUnit{Mesa de soporte}
\SetContactMembers{%
  Equipo CoreX & Soporte & soporte@empresa.com%
}
\SetContactNote{Para soporte técnico o consultas sobre este documento.}

\SetDocType{PLAN DE TRABAJO}
\SetDocTitle{Plan de Trabajo — Período}
\SetDocCode{PLN-XXX-000}
\SetDocArea{Área}
\SetDocAuthor{Responsable}
\SetDocDate{\today}

\begin{document}

\section{Objetivo del período}

\begin{ParrafoEjecutivo}
Describir el objetivo principal del período y los entregables esperados.
\end{ParrafoEjecutivo}

\section{Actividades planificadas}

\begin{table}[H]
  \centering
  \caption{Actividades del período}
  \begin{adjustbox}{max width=\linewidth}
  \begin{tabular}{l l l c l}
    \toprule
    \textbf{Actividad} & \textbf{Responsable} & \textbf{Fecha límite} & \textbf{Prioridad} & \textbf{Estado} \\
    \midrule
    Actividad 1 & Responsable A & 2026-04-30 & Alta    & Pendiente \\
    Actividad 2 & Responsable B & 2026-05-05 & Media   & En curso  \\
    Actividad 3 & Responsable A & 2026-05-10 & Normal  & Pendiente \\
    \bottomrule
  \end{tabular}
  \end{adjustbox}
\end{table}

\section{Recursos requeridos}

\begin{itemize}
  \item Recurso humano: descripción de perfiles necesarios.
  \item Recursos técnicos: herramientas, accesos o equipos.
  \item Tiempo estimado total: X horas.
\end{itemize}

\section{Riesgos identificados}

\begin{table}[H]
  \centering
  \caption{Registro de riesgos}
  \begin{tabular}{p{5cm} c c p{5cm}}
    \toprule
    \textbf{Riesgo} & \textbf{Probabilidad} & \textbf{Impacto} & \textbf{Mitigación} \\
    \midrule
    Descripción del riesgo 1 & Media & Alto & Plan de contingencia. \\
    Descripción del riesgo 2 & Baja  & Medio & Medida preventiva.  \\
    \bottomrule
  \end{tabular}
\end{table}

\begin{ObservationBox}[Dependencias]
  Describir dependencias con otras áreas o sistemas que pueden afectar la ejecución.
\end{ObservationBox}

\SignatureBlock{Nombre Apellido}{Responsable del Plan}

\ContactSignature

\end{document}
